
\documentclass[12pt]{article}
\usepackage[utf8]{inputenc}
\usepackage{amsmath, amssymb}
\usepackage{geometry}
\geometry{margin=1in}
\usepackage{graphicx}
\usepackage{caption}
\usepackage{color}

\title{Practical Example: Albatross Take-Off}
\author{}
\date{}

\begin{document}

\maketitle

\section*{Objective of the Example}

Estimate the forces and moments (torques) acting on an albatross wing during take-off, using a simple mathematical model based on basic physics.

\section*{Problem Statement}

When an albatross takes off, the air exerts a force on its wings. This force depends on:
\begin{itemize}
    \item the airspeed,
    \item the size and shape of the wing,
    \item the angle at which the air hits the wing.
\end{itemize}

We want to estimate:
\begin{enumerate}
    \item How much the air "pushes" the wing upward (lift force),
    \item How much this force makes the wing want to bend (bending moment),
    \item How much this force tries to twist the wing (torsional moment).
\end{enumerate}

\section*{Known Data}

\begin{itemize}
    \item Wing semi-span: $b/2 = 1$ m
    \item Root chord (width at the center): $c_r = 0.2$ m
    \item Chord distribution: $c(z) = c_r \sqrt{1 - 4(z/b)^2}$
    \item Airspeed: $U_\infty = 15$ m/s
    \item Air density: $\rho = 1.225$ kg/m$^3$
    \item Angle of attack: $\alpha = 10^\circ = 0.1745$ rad
    \item Lift coefficient: $C_L = 2\pi \alpha \approx 1.096$
    \item Drag coefficient: $C_D = 0.04$
    \item Elastic axis – aerodynamic center offset: $x_{EA} - x_{AC} = -0.02$ m
\end{itemize}

\section*{Step 1: Forces on the Wing}

The aerodynamic pressure on the wing per unit area is:

\[
F = \frac{1}{2} \rho U^2
\]

This pressure is multiplied by the wing area and the lift coefficient $C_L$ to get the upward lift force:

\[
F_L(z) = \frac{1}{2} \rho U^2 c(z) C_L
\]

We also break this force into two components:
\begin{itemize}
    \item $C_N$: normal to the wing surface.
    \item $C_X$: along the chord of the wing.
\end{itemize}

Calculations:

\[
C_N = C_L \cos\alpha - C_D \sin\alpha \approx 1.071
\]
\[
C_X = C_L \sin\alpha + C_D \cos\alpha \approx 0.229
\]

\section*{Step 2: Bending Moment (Flapwise Moment)}

The further from the root, the larger the moment arm. The bending moment is:

\[
M_F = \int_0^{b/2} \left( \frac{1}{2} \rho U^2 c(z) C_N \right) z \, dz
\]

With an elliptical chord distribution, this becomes:

\[
M_F = \frac{c_r b^2}{24} \rho U^2 C_N
\]

Substituting:

\[
M_F = \frac{0.2 \cdot (2)^2}{24} \cdot 1.225 \cdot 15^2 \cdot 1.071 \approx \boxed{9.84 \, \text{Nm}}
\]

\section*{Step 3: Edgewise Moment}

This is the aerodynamic moment along the flight direction:

\[
M_E = \frac{c_r b^2}{24} \rho U^2 C_X
\]
\[
M_E \approx \frac{0.2 \cdot 4}{24} \cdot 1.225 \cdot 225 \cdot 0.229 \approx \boxed{2.10 \, \text{Nm}}
\]

\section*{Step 4: Torsional Moment}

Due to the offset between the aerodynamic force center and the structural torsion axis:

\[
M_T = \frac{\pi c_r b}{16} \rho U^2 C_R (x_{EA} - x_{AC})
\]

Where:

\[
C_R = \sqrt{C_L^2 + C_D^2} \approx 1.097
\]

Substituting:

\[
M_T = \frac{\pi \cdot 0.2 \cdot 2}{16} \cdot 1.225 \cdot 225 \cdot 1.097 \cdot (-0.02) \approx \boxed{-3.55 \, \text{Nm}}
\]

The negative sign indicates the twist is opposite to the reference direction.

\section*{Conclusion}

During take-off, the albatross wing experiences:

\begin{itemize}
    \item A strong bending moment: 9.84 Nm
    \item A smaller edgewise moment: 2.10 Nm
    \item A significant torsional moment: -3.55 Nm
\end{itemize}

This explains why wings must be very stiff, especially in torsion, to avoid dangerous vibrations or excessive deformation during flight.

\end{document}
